% ! TeX root = ../main.tex


% # ORCA Intro — Working Draft

% ---

% ## Internal Notes / Open Questions

% - **Contributions enumeration**: Some papers do "3 key contributions: (1)... (2)... (3)..." — do
% we want this? Could go at end of P3 or as a standalone mini-para between P3 and P4. Risk:
% formulaic. Alternative: let P3 (system) and P4 (results) carry the contributions implicitly.
% Decision deferred. - **Open source**: Codebase is open source. Currently at end of P3. Don't bury
% it. - **AMR calibration**: How much AMR in P1? Currently one sentence. Could be zero (just "in our
% experience") or could name it. Keep light — §2.2 has the detail. - **Audience**: Systems + HPC.
% Database people are secondary but should find the OLAP framing legible.

% ## Paragraph Roles

% - **P1 — Stakes + Problem**: BSP at scale, efficiency matters, pathologies are context-specific
% and can only be found in the run where they occur. Post-hoc doesn't work. "Why now": scale is
% forcing the issue. - **P2 — Always-on application-aware observability**: Not metrics —
% application-aware. Dapper parallel (requests vs collectives). ML ad-hoc pipelines don't compose.
% The capability model: monitor collectives cheaply, turn on detail when anomalous. Requires
% BSP-aware guarantees. - **P3 — ORCA**: Tree-based overlay. Progression of three capabilities: (1)
% better offline tracing via causal partitioning, (2) streaming in-situ analytics, (3) control plane
% for online intervention. MRNet lineage + modern OLAP primitives. - **P4 — Results + broader
% vision**: Headline eval numbers. Close with the wider claim (1–2 sentences).

% ## Aside: Trimmed Content

% **Storage efficiency numbers** (cut from P4 for density — available if we want them back): -
% Parquet output 2–5× smaller than OTF2, 10–20× smaller than plaintext formats. - Per-record
% overhead remains constant at scale.

% **"Three capabilities in progression"** (from bullets, not explicitly serialized): - The original
% outline had "efficient columnar collection → in-situ SQL analytics → closed-loop control via
% TS2PC" as a three-stage progression. P3 currently covers all three but doesn't frame them as an
% evolution. The paper's §4 intro uses this progression framing. Could reintroduce if we want that
% narrative arc in the intro.

% **HPC/ML convergence** (partially absorbed into P1): - The explicit point that "HPC and ML
% training are converging on the same pathologies but with separate tooling" got compressed. Could
% be made more explicit.

% ## Concerns

% 1. **P2 is long and dense.** It does: metrics insufficient → application-aware → Dapper parallel →
% collectives as BSP primitive → capability model → ML ad-hoc → common substrate + guarantees.
% That's a lot of moves. Might need to split, but the logic is sequential so it may hold as one
% dense paragraph. Defer to formatting pass. 2. **P3 might split** at "ORCA builds on..." —
% capability vs lineage/why-now as separate typeset paragraphs. 3. **"Foundational infrastructure"
% positioning** is implicit, not explicit. Strong enough, or do we want an explicit claim? 4.
% **Contributions enumeration**: Still undecided. 5. **Score-P and Kineto no longer named in P1.**
% P1 now says "capturing massive traces" generically. This is cleaner but loses the concrete
% signaling that we know the landscape. The specifics live in §2 now. Acceptable? 6. **Dapper is
% named but not cited.** Need a citation if we keep the name, or just say "distributed tracing
% infrastructure" without naming it.

% --- ---

% # Prose


The \emph{bulk-synchronous parallel} (BSP)~\cite{Valia1990:BSP} paradigm powers critical
computational workloads, from
scientific simulations to large-scale machine learning training.
These systems are uniquely fragile: at the scale of tens of thousands of processors, lockstep
execution means a single straggler can block cluster-wide progress at every synchronization point.
As ML training pushes BSP toward 100,000 GPUs and scientific workloads grow in tandem, the cost of
undiagnosed performance pathologies has become economically untenable.
% Kong2022: Collie (RDMA) Xuan2026: STAD You2024:GVARP: GVARP Missing: Petrini2003 Acun2016:
% Mitigating Acun2016: Variation Gunaw2018: FailSlow VanS2009: Scalability in AMR
Yet the causes are diverse and highly
context-dependent~\cite{Kong2022:Collie,Xuan2025:STAD,You2024:GVARP,Petri2003:Missing,Acun2016:Mitigating,Acun2016:Variation,Gunaw2018:FailSlow,VanS2009:AMR,Jain2025:AMR,Bhate2020:PerfVar,Chen2025:Reveal,Gangi2024:RDMA,Cui2025:XPUTimer,Jiang2025:TrainCheck,Deng2025:Minder,Dong2025:Aegis,Yao2025:Holmes,Wu2025:GREYHOUND}
---load imbalance, poor overlap, hardware variability---varying with the cluster, workload, scale,
and configuration.
Metrics~\cite{:Prometheus,Schwa2021:LDMS} provide baseline visibility but rarely enough workload
structure to attribute causes.
Tracing~\cite{Shend2006:TAU,Knupf2012:Score-P,2025:PytorchKineto,Devar2024:DFTracer,Boehm2016:Caliper}
can capture the fine-grained patterns needed for diagnosis, but prevailing approaches---exhaustive
collection to disk followed by expensive \emph{extract-transform-load}
(ETL)~\cite{Yildi2023:IOMax,Jain2025:AMR} ---cannot provide feedback at the speed problems demand.
In our study of placement in AMR codes~\cite{Jain2025:AMR}, isolating the root causes of several
performance artifacts took months, while fixes typically took days.
Rapid visibility into production pathologies remains one of the highest-yield paths to efficiency.

We argue for always-on, steerable observability as a production capability for BSP workloads---not
just faster post-hoc analysis, but continuous visibility with the ability to escalate on demand.
The evolution of distributed tracing~\cite{Barha2004:Magpie,Fonse2007:X-Trace,Sigel2010:Dapper}
offers an instructive parallel.
In microservices, tracing evolved to preserve causal relationships between requests---the unit of
work in asynchronous systems---so tail latency is continuously monitored and attributable to
specific services.
% In BSP, causal relationships come from collectives Causal relationships in BSP come from
% collectives: 
In BSP, collectives induce causality: performance anomalies manifest as stragglers that stall
collectives.
And because BSP is a single coordinated workload, the system can unify observation and
intervention---enabling operators to steer collection on demand, test hypotheses live, and act on
feedback, all while the job runs.
Real-time diagnostics are already proving valuable in ML
training~\cite{Yao2025:Holmes,Wu2025:GREYHOUND,Deng2025:Minder,Cui2025:XPUTimer,Dong2025:Aegis,Jiang2025:TrainCheck,Guan2025:PerfTracker}
for problems ranging from straggler localization to training correctness, but each system
implements a bespoke collection and analytics pipeline.
A common observability substrate is needed for BSP workloads, providing low-overhead
on-fabric transport, BSP-aware coordination semantics for telemetry aggregation and control, and
programmable interfaces for analytics (\S\ref{sec:orca:req}).

% We present ORCA (\emph{Observability with Real-time Control and Aggregation}), a tree-based
% overlay network for real-time observability and control of large-scale BSP applications.
We present a \emph{tree-based overlay network} (TBON) for real-time observability and control of
large-scale BSP applications, called
ORCA~(\emph{Observability~with~Real-time~Control~and~Aggregation}).
ORCA treats telemetry as columnar streams, transformed via a SQL-based analytics interface called
\emph{OrcaFlow} en route to \emph{sinks} such as dashboards or parallel filesystems.
A key contribution is the \emph{timestep dataframe} abstraction: building on the insight that
global synchronization divides BSP telemetry into causally independent windows, ORCA transforms
telemetry from a rank-major to a timestep-major representation (\S\ref{sec:orca:des_model}).
An RDMA-based TBON-optimized transport (\S\ref{sec:orca:des_xport}) enables low-overhead collection
of
these dataframes.
They can then be transformed via SQL-based analytics, routed to a dashboard, or persisted as
partitioned Parquet for zero-ETL post-mortems.
Distributed query planning with operator pushdown minimizes data movement across the overlay
(\S\ref{sec:orca:des_insitu}).
A control plane (\S\ref{sec:orca:des_ctl}) provides atomic and timestep-consistent control
semantics,
guaranteeing correctness for interventions such as new analytics, probes, parameter updates.
ORCA combines a modern columnar analytics
stack~\cite{:ApacheArrow,:ApacheParquet,Lamb2024:DataFusion} with RDMA and BSP-aligned data and
control paths to realize low-overhead, real-time, closed-loop feedback.

We evaluate ORCA on a real AMR code at up to 4096 CPU ranks.
% , and progressively study its transport, analytics, and control capabilities.

\begin{itemize}[leftmargin=*]
  \item In conventional tracing workflows, ORCA incurs \textasciitilde{2}\% runtime overhead
        (vs 56--81\% for state-of-the-art tracers), produces up to 42$\times$ smaller
        traces for equivalent telemetry, and enables up to
        6900$\times$ faster analyses than existing trace formats (\S\ref{sec:orca:eval-trace}).
  \item In-situ analytics enable low-overhead needle-in-the-haystack
        workflows---outlier \mpiw{} calls are isolated in real-time at sub-2\% runtime
        overhead, while discarding up to 99.999\% of data at MPI ranks (\S\ref{sec:orca:eval-insitu}).
  \item In an agentic evaluation of its closed-loop capabilities using an LLM
        agent as a proxy for a human operator, ORCA enables localizing a performance
        anomaly---the same bug from our AMR study that took months of manual
        investigation---in 27 minutes, while reducing data volume by up to
        144$\times$ (\S\ref{sec:orca:eval-agent}).
\end{itemize}
More broadly, ORCA demonstrates that declarative SQL analytics can be embedded on the fabric
alongside a live MPI application, enabling streaming analytics workflows beyond observability.
ORCA is open source~\cite{Autho2026:orca-code}.

% We evaluate ORCA on a real AMR code at upto 4096 CPU ranks. Key results from the evaluation are
% summarized below.

% 1. ORCA outperforms conventional tracers in runtime overhead, storage efficiency, and analytics
% performance. 2. In-situ analytics enable needle-in-the-haystack workflows, filtering up to
% 99.999\% of data at MPI ranks while retaining sub-2\% runtime overhead. 3. In an agentic
% evaluation of its closed-loop capabilities using an LLM agent as a proxy for human operator, ORCA
% localizes a performance anomaly --- the same bug from our AMR study that took months of manual
% investigation --- in 27 minutes, while reducing data volume by up to 144x.

% More broadly, ORCA demonstrates that declarative SQL analytics can be embedded on the fabric
% alongside a live MPI application, enabling streaming analytics workflows beyond observability.

% ORCA collects detailed traces at 1–2\% runtime overhead, compared to 56–81\% for state-of-the-art
% tracers whose uncoordinated buffer flushes compound at scale. Post-mortem queries over ORCA's
% Parquet output complete in hundreds of milliseconds—3 to 4 orders of magnitude faster than
% equivalent queries over conventional trace formats, which are bottlenecked by parsing rather than
% computation. In-situ analytics via operator pushdown reduce transmitted data volume by
% 98.5--99.999\% at the MPI tier while maintaining sub-2\% overhead, enabling continuous
% observability that would be impractical with capture-everything tracing. To validate the
% closed-loop workflow, an LLM agent acting as a proxy for a human operator localizes a performance
% anomaly---the same bug from our AMR study that took months to diagnose manually---in 27 minutes.
% More broadly, ORCA demonstrates that declarative SQL analytics can be embedded on the fabric
% alongside a live MPI application, enabling streaming analytics workflows beyond observability.

% The bulk synchronous parallel paradigm powers some of our most critical computational tools,
% ranging from all scientific simulations (weather, climate, fluids, engineering) to large-scale
% distributed machine learning training. Running at scales of up to 100K GPUs, these are some of the
% largest computational workloads out there. Given the resource footprints --- money, energy,
% hardware --- the efficiency of these workloads is of paramount importance.

% We have a language and taxonomy of observability from microservice-based systems that originated
% with Dapper. No such thing for BSP. Production clusters rely on basic metric-level visibility but
% there is no structured solution for deeper application-level visibility. Former can only provide
% basic high-level visibility, but no insight into complex application-level behavior.

% To add to that, a supercomputer is basically the craziest thing computer scientists have built.
% Frequent global synchronization and no appetite for stragglers. Achieved by connecting nodes
% together by exotic fabrics that split the difference between a PCI bus and a network.

% There's a gap between the trace and the view. A trace is write-optimized. The view is a query. The
% gap between the trace and the view is inefficiencies. Spatial and temporal gap. Spatial gap: trace
% is such a large dataset that someone else who does not know the application analyzes it, without
% hardware context. Temporal gap: latency of intervention, challenge of explaining possibly
% transient anomalies.

% How do we optimize the view? Indexing. Streaming indexing. But what dimension to index on? This is
% a basic database strategy: think about the dimension the query will access. The dimension that
% captures the causality intrinsic in the workload. For distributed systems it is the RPC. For
% parallel systems it is the thing that makes them different from distributed: the sync!

% The ideal observability spectrum. The ability to request exactly the needed information from any
% layer, aggregate it as fast as
% possible, with zero workload interference.
